\scanpage{3}
\mspage{1}
Catégories \unsure{localisées} et \unsure{théorie|th.} homotopique

\uline{1}. La catégorie des chemins.

Soit $\Phi$ la catégorie suivante
\[
\left\{\begin{aligned}
\operatorname{Ob}\Phi &\simeq \mathbb Z &&\text{On note }\{\alpha\}\ \inserted{\text{ou simplement }\alpha}\ \text{l'objet qui correspond à }\alpha\in\mathbb Z\\
\operatorname{Fl}\Phi &\simeq \inserted{\underbrace{\operatorname{Ob}\Phi}_{\text{identités}}\amalg}\mathbb Z &&\text{On note }\unsure{\underline{u}|u}_{\alpha}\ \text{la flèche qui \unsure{correspond|corresp.}}\\
&&&\text{à l'entier }\alpha\in Z\\
s(u_\alpha)&=\alpha\ \text{ si }\alpha\text{ pair},\ \alpha+1 &&\text{si }\alpha\text{ impair}\\
b(u_\alpha)&=\alpha+1\ \text{ si }\alpha\text{ pair},\ \alpha-1 &&\text{si }\alpha\text{ impair}
\end{aligned}\right.
\]
Il n'y a pas à définir la composition, car si deux flèches sont composables, \unsure{elles sont} l'identité. En fait, $\Phi$ est une catégorie \uline{ordonnée}, $\{\alpha\}\leftarrow\{\beta\}$ ssi $\alpha$ pair, $\beta=\alpha+1$.

On \deleted{\unsure{pose} $\unsure{=}$} note, pour $n\in\mathbb N$
\[
\Phi_n\subset\Phi\qquad \operatorname{Ob}\Phi_n=\{\alpha\in\operatorname{Ob}\Phi\mid\deleted{|\alpha|\leq n}\inserted{-n\leq\alpha\leq+n}\}
\]
la sous-catégorie pleine de $\Phi$ dont les objets sont les $\{\alpha\}$ avec $-n\leq\alpha\leq+n$. Ainsi
\diagram{\begin{tikzcd}[ampersand replacement=\&,column sep=1.1em]
\Phi_1=\bigl\{\ \& \underset{-1}{\bullet} \& \underset{0}{\bullet}\arrow[l]\arrow[r] \& \underset{1}{\bullet} \& \bigr\}\,,
\end{tikzcd}}
la duale est
\diagram{\begin{tikzcd}[ampersand replacement=\&,column sep=1em]
\Psi_1=\Phi_1^{o}=\bigl\{ \& {-1}\arrow[r] \& 0 \& 1\arrow[l] \& \bigr\}
\end{tikzcd}}

On peut considérer $\Phi$ comme somme \unsure{amalgamée} d'une suite ordonnée \inserted{infinie (des deux côtés)} de copies de $\Phi_1$, ou aussi
\diagram{\begin{tikzcd}[ampersand replacement=\&,column sep=.48em,scale=.76,transform shape,row sep=.25em]
\cdots \& \& \smash{\overbrace{\hspace{3.8em}}^{\Phi_1}} \& \& \smash{\overbrace{\hspace{3.8em}}^{\Phi_{\unsure{1}}}} \& \& \smash{\overbrace{\hspace{3.8em}}^{\Phi_1}} \& \& \smash{\overbrace{\hspace{3.8em}}^{\Phi_1}} \& \& \smash{\overbrace{\hspace{3.8em}}^{\Phi_1}} \& \& \cdots\\
\& {} \& \underset{-4}{\cdot}\arrow[l]\arrow[r] \& \underset{-3}{\cdot} \& \underset{-2}{\cdot}\arrow[l]\arrow[r] \& \underset{-1}{\cdot} \& \underset{0}{\bullet}\arrow[l]\arrow[r] \& \underset{1}{\cdot} \& \underset{2}{\cdot}\arrow[l]\arrow[r] \& \underset{3}{\cdot} \& \underset{4}{\cdot}\arrow[l]\arrow[r] \& {}\\
\cdots \& \& \& \smash{\underbrace{\hspace{3.8em}}_{\Psi_1}} \& \& \smash{\underbrace{\hspace{3.8em}}_{\Psi_1}} \& \& \smash{\underbrace{\hspace{3.8em}}_{\Psi_1}} \& \& \smash{\underbrace{\hspace{3.8em}}_{\Psi_1}} \& \& \cdots
\end{tikzcd}}
d'une suite similaire de copies de $\Psi_1$. Les $\Phi_n$ \inserted{pour $n$ impair, sont des \unsure{$=$} sommes de copies de $\Phi_1$ mises bout à bout et les \deleted{\illegible{1}} $\Phi_n$ pour $n$ pair $>0$, \unsure{de} \inserted{\unsure{sommes} de copies de $\Psi_1$}} % slanted
forment \unsure{un|une} \deleted{\unsure{suite infinie}} \inserted{système inductif \unsure{fil.}} avec foncteurs de transition fidèles, dont $\Phi$ est la limite inductive
\[
\Phi=\varinjlim_n\Phi_n.
\]

Mais pour la suite, on va faire des $\Phi_n$ un système \uline{projectif}, en définissant pour tout $n\leq m$ une \uline{rétraction} $\underline{\unsure{f|p|\rho}}_{m,n}$ de $\Phi_m$ sur $\Phi_n$, par

% archive sheet number: 2
